% ---------------------------------------------------------------------
%  CHƯƠNG 2 — CƠ SỞ LÝ THUYẾT VÀ CÔNG NGHỆ
% ---------------------------------------------------------------------
\chapter{Cơ sở lý thuyết và công nghệ}
\label{ch:coso}

\begin{muctieu}
Chương này trình bày các cơ sở lý thuyết và công nghệ thực tế được sử dụng
trong hệ thống trợ lý thị giác AI. Nội dung tập trung vào AI đa phương thức,
thị giác máy tính, OCR, nhận dạng và tổng hợp giọng nói, giao diện Streamlit
và cơ chế lưu lịch sử/cấu hình phiên làm việc.
\end{muctieu}

\section{Kiến thức nền tảng}

\subsection{Trí tuệ nhân tạo và AI đa phương thức}

Trí tuệ nhân tạo \en{(Artificial Intelligence -- AI)} là lĩnh vực nghiên cứu
các phương pháp giúp máy tính thực hiện những nhiệm vụ thường cần khả năng
nhận thức và xử lý thông tin của con người. Trong đồ án, AI được sử dụng để
phân tích hình ảnh, nhận biết vật thể, bối cảnh, nguy cơ, chữ và mệnh giá tiền,
sau đó tạo phản hồi bằng tiếng Việt.

Hệ thống có tính chất đa phương thức vì tiếp nhận hình ảnh, văn bản và bản ghi
âm thanh đã được chuyển thành văn bản, đồng thời tạo ra văn bản và âm thanh phản
hồi. Thành phần Vision sử dụng Gemini thông qua Google GenAI SDK để phân tích
hình ảnh; API hỗ trợ đầu ra JSON có cấu trúc
~\cite{GoogleGeminiVision2026,GoogleGeminiStructuredOutput2026,GoogleGenAIPython2026}.
Trong source, Gemini được yêu cầu trả JSON qua MIME type; chương trình chuẩn hoá
các trường nhận được trước khi hiển thị. Thành phần giọng nói đảm nhiệm STT và TTS.

\subsection{Thị giác máy tính và Vision Language Model}

Thị giác máy tính \en{(Computer Vision)} cho phép máy tính trích xuất thông tin
có ý nghĩa từ hình ảnh. Trong hệ thống, ảnh được tải lên hoặc chụp từ camera,
sau đó được chuyển thành dữ liệu phù hợp để gửi đến Gemini Vision
~\cite{PillowImageModule2026}.

Mô hình Vision được yêu cầu trả về các trường có cấu trúc như mô tả, độ tin
cậy, bối cảnh, vật thể, nguy cơ, hướng dẫn di chuyển, văn bản nhận diện và
mệnh giá tiền. Cách trả kết quả dạng JSON giúp chương trình Python có thể xử
lý và hiển thị từng thành phần một cách ổn định.

Ở mức nền tảng, Transformer và Vision Transformer là các hướng nghiên cứu áp
dụng attention cho dữ liệu chuỗi và ảnh; CLIP là ví dụ về học liên kết ảnh với
ngôn ngữ tự nhiên. Đây là tài liệu tham khảo lý thuyết, không phải các mô hình
được gọi trực tiếp trong source của đồ án
~\cite{Vaswani2017Attention,Dosovitskiy2021ViT,Radford2021CLIP}.

Trong phiên bản hiện tại, việc chuẩn hoá ảnh trước khi gửi giúp giới hạn kích
thước dữ liệu truyền; kết quả JSON giúp giao diện truy cập các trường theo
tên thay vì phân tích văn bản tự do. Đây là lựa chọn về luồng dữ liệu, không
phải cam kết về thời gian phản hồi cố định.

\subsection{Nhận dạng ký tự quang học (OCR)}

OCR \en{(Optical Character Recognition)} là kỹ thuật chuyển nội dung chữ
trong hình ảnh thành văn bản. Đồ án sử dụng Tesseract OCR chạy cục bộ cho
chế độ ``Đọc chữ trong ảnh (OCR)''. Đây là một nhánh xử lý riêng, không giả
lập OCR bằng Gemini.

Trước khi nhận dạng, ảnh được chuyển sang RGB, phóng to khi cần, tăng tương
phản, tăng độ sắc nét và tạo thêm biến thể nhị phân. Tesseract được chạy với
các chế độ phân đoạn phù hợp; dữ liệu tiếng Việt \texttt{vie} và tiếng Anh \texttt{eng} được
ưu tiên khi có sẵn. Kết quả được lọc theo độ tin cậy và loại bỏ các chuỗi
nhiễu trước khi hiển thị. Tesseract là bộ máy OCR mã nguồn mở; tài liệu kỹ
thuật của Tesseract mô tả quy trình nhận dạng văn bản và các thành phần xử lý
ảnh~\cite{Smith2007Tesseract,TesseractDocumentation2026}.

\subsection{Nhận dạng giọng nói (Speech-to-Text)}

Speech-to-Text \en{(STT)} chuyển giọng nói thành văn bản để hệ thống có thể
xử lý câu hỏi của người dùng. Phiên bản hiện tại dùng thư viện
\texttt{SpeechRecognition} gọi Google Web Speech để nhận dạng tiếng Việt.
Luồng này cần kết nối Internet; âm thanh chỉ được gửi sau khi người dùng chủ
động ghi âm và xác nhận xử lý. Source không tích hợp Faster-Whisper hoặc STT
chạy cục bộ~\cite{SpeechRecognitionPackage2026}.

\subsection{Tổng hợp giọng nói (Text-to-Speech)}

Text-to-Speech \en{(TTS)} chuyển phản hồi văn bản thành âm thanh. Hệ thống
ưu tiên Edge TTS với giọng tiếng Việt \texttt{vi-VN-HoaiMyNeural}~\cite{EdgeTTS2026}. Nếu Edge TTS gặp
lỗi, ứng dụng chuyển sang Windows SAPI5 thông qua \texttt{pyttsx3}. File âm thanh
được tạo tạm thời và phát bằng thành phần audio của Streamlit.

\subsection{Khả năng tiếp cận cho người khiếm thị}

Hệ thống hướng tới việc giảm phụ thuộc vào quan sát màn hình bằng cách cho
phép người dùng tải/chụp ảnh, đặt câu hỏi bằng văn bản hoặc giọng nói và
nghe kết quả. Các chức năng hỗ trợ gồm mô tả ảnh, đọc chữ, nhận diện tiền,
phát hiện nguy cơ, hỗ trợ di chuyển và đọc kết quả bằng giọng nói. Các tiêu
chí về focus bàn phím, tương phản và giảm chuyển động có liên quan đến hướng
dẫn tiếp cận nội dung web WCAG~\cite{W3CWCAG22}. Hướng dẫn WAI về văn bản thay
thế cũng phân biệt ảnh thông tin, ảnh trang trí và ảnh chức năng
~\cite{W3CImagesTutorial2026}.

\section{Công cụ và công nghệ}

Bảng dưới đây được tổng hợp trực tiếp từ \texttt{requirements.txt}, \texttt{.env.example}
và cấu trúc mã nguồn của phiên bản VisionAI Pro.

\begin{table}[H]
\centering
\caption{Các công cụ, thư viện và công nghệ được sử dụng trong đồ án.}
\label{tab:congcu}
\begin{tabular*}{\textwidth}{@{\extracolsep{\fill}}>{\bfseries\color{daunavy}}p{3.1cm} >{\color{daumaroon}}p{3.1cm} p{7.0cm}}
\toprule
\rowcolor{gray100}
Thành phần & Phiên bản khai báo tối thiểu & Vai trò trong hệ thống \\
\midrule
Python & Python 3.x & Ngôn ngữ lập trình chính cho giao diện, xử lý ảnh, AI và giọng nói. \\
Streamlit & $\ge$ 1.50 & Framework xây dựng giao diện web chạy cục bộ. \\
Google GenAI SDK & $\ge$ 1.30.0 & Kết nối Gemini để phân tích ảnh và hỏi đáp theo ảnh. \\
Gemini Vision & ID mặc định trong source: \texttt{gemini-3.8-flash} & Đã kiểm tra kết nối bằng \texttt{gemini-3.5-flash-lite}; các model còn lại và độ chính xác chưa được đánh giá. \\
Pillow & $\ge$ 10.0 & Đọc và xử lý ảnh đầu vào. \\
Tesseract + pytesseract & pytesseract $\ge$ 0.3.13 & OCR cục bộ, ưu tiên dữ liệu tiếng Việt và tiếng Anh. \\
SpeechRecognition & $\ge$ 3.10 & STT tiếng Việt thông qua Google Web Speech; cần Internet. \\
Edge TTS & $\ge$ 7.0 & TTS tiếng Việt, ưu tiên giọng \texttt{vi-VN-HoaiMyNeural}. \\
pyttsx3 & $\ge$ 2.98 & TTS dự phòng bằng Windows SAPI5. \\
python-dotenv & $\ge$ 1.0 & Đọc biến môi trường và API key từ file \texttt{.env}. \\
JSON & Có sẵn trong Python & Lưu lịch sử phân tích tại \texttt{data/history.json}, tối đa 100 bản ghi. \\
\bottomrule
\end{tabular*}
\end{table}

Bảng~\ref{tab:congcu} cho biết thư viện nào đảm nhiệm từng công đoạn, nhờ đó
việc bảo trì và kiểm thử có thể tập trung vào từng thành phần. Tách riêng OCR,
STT và TTS cũng giúp chọn dịch vụ phù hợp cho từng loại đầu vào; riêng OCR
chạy cục bộ nên không phải tải ảnh lên Gemini khi người dùng chỉ cần đọc chữ.
Ứng dụng giao diện được xây dựng trên Streamlit~\cite{StreamlitDocumentation2026}.
Các thao tác lấy ảnh từ webcam và ghi âm dựa trên widget
\texttt{st.camera\_input} và \texttt{st.audio\_input} của Streamlit
~\cite{StreamlitCameraInput2026,StreamlitAudioInput2026}.

Không sử dụng cơ sở dữ liệu quan hệ trong phiên bản source hiện tại. Lịch sử
được lưu bằng tệp JSON cục bộ; cấu hình giọng nói, mức chi tiết và trạng thái
hội thoại được quản lý bằng \texttt{st.session\_state} trong phiên Streamlit.

\section{Mô hình AI và cơ chế dự phòng}

Danh sách ID model Gemini được cấu hình theo thứ tự trong biến
\texttt{GEMINI\_MODELS}; giá trị mặc định trong source là
\texttt{gemini-3.8-flash,gemini-3.6-flash,gemini-3.5-flash-lite}. Ứng dụng lần lượt
thử các model khi một yêu cầu gặp lỗi, dừng khi API từ chối khoá/quyền truy cập
(401/403), và chờ ngắn trước khi chuyển model khi gặp giới hạn tốc độ (429).
Danh sách dự phòng không khắc phục được khoá sai, hết hạn mức toàn tài khoản
hoặc lỗi mạng; khả năng truy cập đã được xác nhận với
\texttt{gemini-3.5-flash-lite}, còn các ID khác trong danh sách chưa được kiểm
tra bằng yêu cầu API thật. Cài đặt có nút gửi một yêu cầu nhỏ để kiểm tra
kết nối Gemini.

Yêu cầu Vision sử dụng JSON response mode. Source không khai báo công cụ
function calling; kết quả được chuẩn hoá về một cấu trúc thống nhất trước khi
giao diện hiển thị.

\section{Định hướng lựa chọn công nghệ}

Source tách các tác vụ theo khả năng và phụ thuộc của từng thành phần: Gemini
xử lý mô tả ảnh và hỏi đáp; Tesseract cung cấp OCR cục bộ; SpeechRecognition
chuyển bản ghi thành văn bản qua dịch vụ trực tuyến; Edge TTS và Windows SAPI5
tạo âm thanh phản hồi. Cách phân tách này cho phép dùng OCR khi không cần gọi
Gemini, nhưng không đồng nghĩa toàn bộ ứng dụng chạy ngoại tuyến. Đây là mô tả
kiến trúc source, không phải kết quả so sánh định lượng với hệ thống khác.

\section{Tóm tắt chương}

Chương này đã xác định các công nghệ thực tế của hệ thống gồm Python,
Streamlit, Gemini Vision thông qua Google GenAI SDK, Pillow, Tesseract OCR,
SpeechRecognition/Google Web Speech, Edge TTS, pyttsx3, python-dotenv và JSON.
Các thành phần này tạo thành chuỗi xử lý đa phương thức từ ảnh/giọng nói đến
phân tích AI, văn bản và âm thanh phản hồi.
